\sgaStableTarget{sga6.e09.expose}\section*{第 IX 讲. $K_\bullet$ 群的若干计算}
\begin{center}
P. Berthelot 著
\end{center}

本讲没有对应的口头报告, 后续各讲也不使用它. 其中计算了各种纤维化的 $K_\bullet$ 群: 向量丛、射影丛、线性群下的挠子等. 本讲始终假定基概形是诺特的, 并用 $K_\bullet(X)$ 表示概形 $X$ 上凝聚层的 Grothendieck 群. 回顾一下, 若 $X$ 是诺特概形, 则从凝聚 $\mathcal O_X$ 模范畴的 $K$ 群到 $X$ 上伪凝聚复形范畴的 $K$ 群的同态是同构 (\hyperlink{sga6.e04.para.2-4}{IV 2.4}).

这里的结果与 Chow 环的相应结果类似, 后者可见 Grothendieck 在 1958 年 Chevalley 讨论班中的报告. 还可以问, 对 $K_\bullet$ 按支集维数定义的滤过所对应的分次群 (参见 X), 在多大程度上能得到同类命题. 即便对于\hyperlink{sga6.e09.prop.1-1}{命题 1.1} 所述的同伦正合列, 答案似乎也尚不清楚; 而这一性质是以下计算中使用最多的性质之一.

\sgaStableTarget{sga6.e09.sec.1}\subsection*{1. 向量丛}

\sgaStableTarget{sga6.e09.prop.1-1}\textbf{命题 1.1 (同伦正合列).} 设 $X$ 为诺特概形, $X'$ 为 $X$ 的闭子概形, $U$ 为 $X'$ 在 $X$ 中的开补集, $i$ 为 $X'$ 到 $X$ 的浸入, $j$ 为 $U$ 到 $X$ 的浸入. 则序列
\[
 K_\bullet(X')\xrightarrow{i_*}K_\bullet(X)\xrightarrow{j^*}K_\bullet(U)\longrightarrow 0
\]
正合.

若 $F'$ 为凝聚 $\mathcal O_{X'}$ 模, 则 $i_*(F')$ 是支集包含于 $X'$ 的 $\mathcal O_X$ 模; 因而它在 $U$ 上的限制为零, 且 $j^*i_*(\cl(F'))=0$. 由于形如 $\cl(F')$ 的元素生成 $K_\bullet(X')$, 故 $j^*i_*=0$. 因此得到同态
\[
 \rho:K_\bullet(X)/\operatorname{Im}K_\bullet(X')\longrightarrow K_\bullet(U).
\]
我们来定义 $\rho$ 的逆同态. 设 $G$ 为凝聚 $\mathcal O_U$ 模. 因为 $X$ 是诺特的, 存在一个凝聚 $\mathcal O_X$ 模 $F$, 其在 $U$ 上的限制同构于 $G$. 若 $F'$ 是另一个在 $U$ 上的限制同构于 $G$ 的凝聚 $\mathcal O_X$ 模, 则在 $K_\bullet(X)$ 中有
\[
 \cl(F)-\cl(F')\in \operatorname{Im}K_\bullet(X').
\]
为此, 利用对角映射将 $G$ 看作 $j^*(F\oplus F')$ 的子层. 于是 $G$ 可延拓为 $F\oplus F'$ 的子层 $F''$, 而问题归结为证明
\[
 \cl(F)-\cl(F'')\in \operatorname{Im}K_\bullet(X')
\]
以及对 $F'$ 的类似结论. 设 $H$ 和 $K$ 分别为复合同态
\[
 F''\longrightarrow F\oplus F'\longrightarrow F.
\]
的核与余核. 此同态在 $U$ 上的限制为 $G$ 的恒等映射, 因而 $H$ 和 $K$ 的支集包含于 $X'$. 若 $I$ 是 $X'$ 在 $X$ 中的理想, 则存在整数 $m$ 和 $n$, 使得 $I^mH=0$ 且 $I^nK=0$ (EGA I 5.3.4). 因此可将 $H$ 和 $K$ 逐层分解为被 $I$ 零化的模的扩张, 即 $\mathcal O_{X'}$ 模的扩张, 于是
\[
 \cl(H)\in\operatorname{Im}K_\bullet(X'),\qquad \cl(K)\in\operatorname{Im}K_\bullet(X'),
\]
这就证明了断言. 因此, 给每个凝聚 $\mathcal O_U$ 模对应其在 $X$ 上某个延拓的类, 便定义了凝聚 $\mathcal O_U$ 模范畴上取值于 $K_\bullet(X)/\operatorname{Im}K_\bullet(X')$ 的函数.

此函数是可加的. 事实上, 设
\[
0\longrightarrow G'\longrightarrow G\longrightarrow G''\longrightarrow 0
\]
是凝聚 $\mathcal O_U$ 模的正合列. 设 $F$ 为 $G$ 到 $X$ 上的延拓. 存在 $F$ 的子模 $F'$ 延拓 $G'$, 而 $F/F'$ 延拓 $G''$. 因此此函数经由 $K_\bullet(U)$ 分解, 定义出所需的同态 $K_\bullet(U)\to K_\bullet(X)/\operatorname{Im}K_\bullet(X')$. 立即可验证它是 $\rho$ 的逆.

\sgaStableTarget{sga6.e09.prop.1-2}\textbf{命题 1.2.} 设 $X$ 和 $S$ 为诺特概形, $f:X\to S$ 为平坦态射. 给定 $K^\circ(X)$ 中的一族元素 $(\xi_i)_{i\in I}$, 并考虑同态
\[
 \varphi:K_\bullet(S)^{(I)}\longrightarrow K_\bullet(X)
\]
定义为
\[
 \varphi((x_i)_{i\in I})=\sum_{i\in I}\xi_i\,f^*(x_i),
\]
其中 $f^*$ 是\hyperlink{sga6.e04.para.2-12}{IV 2.12} 中定义的从 $K_\bullet(S)$ 到 $K_\bullet(X)$ 的逆像同态, 并按\hyperlink{sga6.e04.para.2-10}{IV 2.10} 将 $K_\bullet(X)$ 看作 $K^\circ(X)$ 模. 假定对每个点 $s\in S$, 各 $\xi_i$ 经同态
\[
 K^\circ(X)\longrightarrow K^\circ(X_s)\longrightarrow K_\bullet(X_s)
\]
在 $K_\bullet(X_s)$ 中的像生成阿贝尔群 $K_\bullet(X_s)$. 则 $\varphi$ 为满射.

设 $S'$ 是 $S$ 的闭子概形, $U$ 为 $S'$ 在 $S$ 中的开补集. 以 $X_U$ 表示 $X$ 在 $U$ 上的限制, 以 $X_{S'}$ 表示 $X$ 在 $S'$ 上的限制. 仍以 $\xi_i$ 表示各 $\xi_i$ 在 $K^\circ(X_{S'})$ 和 $K^\circ(X_U)$ 中的像, 并以
\[
 \varphi_{S'}:K_\bullet(S')^{(I)}\to K_\bullet(X_{S'}),\qquad
 \varphi_U:K_\bullet(U)^{(I)}\to K_\bullet(X_U)
\]
表示仿照 $\varphi$ 定义的同态.

考虑图表
\sgaStableTarget{sga6.e09.eq.1-2-1}\[
\begin{tikzcd}[column sep=tiny]
K_\bullet(S')^{(I)} \arrow[r,"i_*"] \arrow[d,"\varphi_{S'}"'] &
K_\bullet(S)^{(I)} \arrow[r,"j^*"] \arrow[d,"\varphi"] &
K_\bullet(U)^{(I)} \arrow[r] \arrow[d,"\varphi_U"] & 0\\
K_\bullet(X_{S'}) \arrow[r,"u_*"] & K_\bullet(X) \arrow[r,"v^*"] & K_\bullet(X_U) \arrow[r] & 0,
\end{tikzcd}\tag{1.2.1}
\]
其中 $i$ 和 $j$ 分别为 $S'$ 和 $U$ 到 $S$ 的浸入, $u$ 和 $v$ 分别为 $X_{S'}$ 和 $X_U$ 到 $X$ 的浸入. 我们证明: 若 $\varphi_{S'}$ 和 $\varphi_U$ 为满射, 则 $\varphi$ 也是满射.

由 1.1, 两个横行均正合. 因此, 为证 $\varphi$ 为满射, 只需证明两个方块交换. 对 $x_i\in K_\bullet(S')$, 有
\[
 \varphi(i_*(x_i))=\sum_i \xi_i f^*i_*(x_i),
\]
另一方面,
\[
 u_*\varphi_{S'}((x_i))=u_*\left(\sum_i\xi_i u^*f^*(x_i)\right)=\sum_i \xi_i u_*(u^*f^*(x_i)),
\]
最后的等式来自投影公式\hyperlink{sga6.e04.eq.2-11-1-2}{IV 2.11.1.2}. 因此, 等式 $\varphi i_*=u_*\varphi_{S'}$ 来自 $f^*i_*=u_*u^*f^*$; 后者由\hyperlink{sga6.e04.prop.3-1-1}{IV 3.1.1} 成立, 因为 $i$ 固有且 $f$ 平坦. 同样, 对 $(x_i)\in K_\bullet(S)^{(I)}$, 有
\[
 \varphi_U(j^*(x_i))=\sum_i\xi_i f_U^*j^*(x_i),
\]
及
\[
 v^*\varphi((x_i))=v^*\left(\sum_i\xi_i f^*(x_i)\right)=\sum_i\xi_i v^*f^*(x_i)=\sum_i\xi_i f_U^*j^*(x_i),
\]
中间的等式即\hyperlink{sga6.e04.eq.2-12-2}{IV 2.12.2}.

利用这个结果, 可对 $S$ 的闭子概形集作诺特归纳来证明\hyperlink{sga6.e09.prop.1-2}{1.2}. 若假定对 $S$ 的每个不同于 $S$ 的闭子概形 $S'$, $\varphi_{S'}$ 均为满射, 则上述论证说明, 在必要时将 $S$ 换成某个不可约仿射开集之后, 可归结为不可约仿射基概形上的\hyperlink{sga6.e09.prop.1-2}{1.2}. 此外, 图表
\[
\begin{tikzcd}
K_\bullet(S_{\mathrm{red}})^{(I)} \arrow[r] \arrow[d,"\varphi_{S_{\mathrm{red}}}"'] & K_\bullet(S)^{(I)} \arrow[d,"\varphi"]\\
K_\bullet(X_{S_{\mathrm{red}}}) \arrow[r] & K_\bullet(X)
\end{tikzcd}
\]
中, $K_\bullet(X_{S_{\mathrm{red}}})\to K_\bullet(X)$ 由\hyperlink{sga6.e09.prop.1-1}{1.1} 为满射, 因而还可假定 $S$ 约化.

这样, 基概形为整仿射概形. 置 $S=\Spec(A)$, 以 $K$ 表示 $A$ 的分式域. 若 $s$ 为 $S$ 的泛点, 则 $\mathcal O_{S,s}=K$. 取 $y\in K_\bullet(X)$. 按假设, $K_\bullet(X_s)$ 由各 $\xi_i$ 的像生成; 若 $y_s$ 为 $y$ 在 $K_\bullet(X_s)$ 中的像, 则可找到整数 $n_i$, 使得
\[
 y_s=\sum_i \xi_i n_i.
\]
另一方面, $K=\varinjlim_f A_f$, 其中 $f$ 遍历 $A$ 的元素. 因而 $\Spec(K)=\varprojlim_f \Spec(A_f)$, 且由 EGA IV 8.2.3 和 8.2.5, 有
\[
 X_s=X\times_S\Spec(K)=\varprojlim_f X\times_S\Spec(A_f)
\]
因为这些概形是诺特的, 且过渡态射仿射而平坦, 由\hyperlink{sga6.e04.para.3-2-3}{IV 3.2.3}, 有
\[
 K_\bullet(X_s)=\varinjlim_f K_\bullet(X\times_S\Spec(A_f)).
\]
由于 $K_\bullet$ 上的逆像同态对于 $K^\circ$ 线性 (\hyperlink{sga6.e04.eq.2-12-2}{IV 2.12.2}), 存在 $S$ 的非空开集 $U=\Spec(A_f)$, 使得在 $K_\bullet(X_U)$ 中有
\[
 y_U=\sum_i \xi_i n_i.
\]
设 $S'$ 是支集为 $S-U$ 的 $S$ 的闭子概形. 由归纳假设, $\varphi_{S'}$ 为满射, 因此图表\hyperlink{sga6.e09.eq.1-2-1}{(1.2.1)} 说明
\[
 y-
\sum_i\xi_i n_i\in \operatorname{Im}(\varphi).
\]
由于 $n_i=f^*(n_i)$, $\varphi$ 为满射.

\sgaStableTarget{sga6.e09.cor.1-3}\textbf{推论 1.3.} 设 $f:X\to S$ 为诺特概形间的平坦态射. 假定对 $S$ 的每个点 $s$, 同态
\[
 K^\circ(X)\longrightarrow K_\bullet(X_s)
\]
为满射. 则由 $f^*$ 导出的同态
\[
  K_\bullet(S)\otimes_{K^\circ(S)}K^\circ(X)\longrightarrow K_\bullet(X),
\]
为满射.

取 $(\xi_i)$ 为 $K^\circ(X)$ 的全部元素组成的族. 于是\hyperlink{sga6.e09.prop.1-2}{1.2} 的假设满足, 因而
\[
 \varphi:K_\bullet(S)^{(K^\circ(X))}\longrightarrow K_\bullet(X)
\]
为满射. 按 $\varphi$ 的定义, 这等价于上述张量积同态为满射.

\sgaStableTarget{sga6.e09.cor.1-4}\textbf{推论 1.4.} 设 $f:X\to S$ 为诺特概形间的平坦态射. 假定对 $S$ 的每个点 $s$, 典范同态 $\mathbb Z\to K_\bullet(X_s)$ 均为同构. 则
\[
 f^*:K_\bullet(S)\longrightarrow K_\bullet(X)
\]
为满射. 若进一步 $f$ 有一个有限 Tor 维数的截面 $g$, 则 $f^*$ 为同构.

取 $(\xi_i)$ 为仅含 $1\in K^\circ(X)$ 的族. 则 $\varphi=f^*$; 由于\hyperlink{sga6.e09.prop.1-2}{1.2} 的条件满足, $f^*$ 为满射. 若 $f$ 有有限 Tor 维数的截面 $g$, 则可定义同态
\[
 g^*:K_\bullet(X)\longrightarrow K_\bullet(S)
\]
(\hyperlink{sga6.e04.para.2-12}{IV 2.12}). 由\hyperlink{sga6.e04.eq.2-12-1-1}{IV 2.12.1.1}, $g^*f^*=(fg)^*=\operatorname{Id}$, 因此 $f^*$ 为单射, 从而为同构.

\sgaStableTarget{sga6.e09.rem.1-5}\textbf{注 1.5.} $g$ 具有有限 Tor 维数的充要条件是 $X$ 的纤维在 $g(S)$ 的各点正则, 即 $f$ 在 $g(S)$ 的各点正则; 若假定为有限表示, 则等价于 $f$ 在这些点光滑.

由于这些性质在 $X$ 和 $S$ 上是局部的, 可假定二者仿射: $S=\Spec(A)$, $X=\Spec(B)$. 则有
\[
 \operatorname{tor.dim}_B(A)=\operatorname{dim.proj}_B(A)=\sup_{s\in S}\operatorname{dim.proj}_{B\otimes_A k(s)}(k(s)),
\]
这是 EGA IV 12.3.2 的结论, 其中 $B$ 的有限表示假设可换成诺特性假设. 而 $k(s)$ 同构于 $\mathcal O_{X_s,g(s)}$ 的剩余域, 且由 EGA $0_{\mathrm{IV}}$ 17.2.5, 有
\[
 \operatorname{dim.proj}_{B\otimes_A k(s)}(k(s))=
 \operatorname{dim.proj}_{\mathcal O_{X_s,g(s)}}(k(s)).
\]
由 EGA $0_{\mathrm{IV}}$ 17.2.11, 最后这个维数有限当且仅当 $\mathcal O_{X_s,g(s)}$ 的上同调维数有限, 即该环正则.

\sgaStableTarget{sga6.e09.prop.1-6}\textbf{命题 1.6.} 设 $S$ 为诺特概形, $E$ 为有限型局部自由 $\mathcal O_S$ 模, $X$ 为 $E$ 所伴的向量丛, $f$ 为 $X$ 的结构态射. 则
\[
 f^*:K_\bullet(S)\longrightarrow K_\bullet(X)
\]
为同构.

因为 $f$ 光滑且有截面, 即零截面, 为证明\hyperlink{sga6.e09.prop.1-6}{1.6}, 由\hyperlink{sga6.e09.cor.1-4}{1.4} 只需证明对每个 $s\in S$ 有 $K_\bullet(X_s)\simeq \mathbb Z$. 因此归结为基概形为域的情形. 此时 $E$ 是秩一子模的直和, 可将 $f$ 分解为一串秩一纤维化, 从而归结为 $X$ 的秩为一的情形, 即 $X=\Spec(k[T])$. 命题于是由下一引理得到.

\sgaStableTarget{sga6.e09.lem.1-7}\textbf{引理 1.7.} 设 $A$ 为主理想整环, $X=\Spec(A)$. 则典范同态
\[
 \mathbb Z\longrightarrow K_\bullet(X)
\]
为同构.

因为 $X$ 正则, 同态 $K^\circ(X)\to K_\bullet(X)$ 为同构. 因而 $K_\bullet(X)$ 由有限型局部自由 $\mathcal O_X$ 模的类生成, 即由有限型投射 $A$ 模的类生成. 由于 $A$ 为主理想整环, 每个有限型投射模均自由, 由此得到满射性. 单射性来自 $K^\circ(X)$ 具有 $\mathbb Z$ 增广这一事实.

\sgaStableTarget{sga6.e09.prop.1-8}\textbf{命题 1.8.} 设 $S$ 为诺特概形, $E$ 为秩 $r$ 的局部自由 $\mathcal O_S$ 模, $X$ 为 $E$ 所伴的向量丛, $f$ 为 $X$ 的结构态射, $g$ 为 $X$ 的零截面. 则复合同态
\[
 K_\bullet(S)\xrightarrow{g_*}K_\bullet(X)\xrightarrow{g^*}K_\bullet(S)
\]
等于乘以 $K^\circ(S)$ 中的元素 $\lambda_{-1}(E)$, 其中 $g^*$ 按\hyperlink{sga6.e04.eq.2-12-1}{IV 2.12.1} 及\hyperlink{sga6.e09.rem.1-5}{1.5} 定义, 用 $E$ 也表示 $E$ 在 $K^\circ(S)$ 中的类, 并置
\[
 \lambda_{-1}(E)=\sum_{i=0}^r(-1)^i\lambda^i(E).
\]

此命题是\hyperlink{sga6.e07.para.2-6}{VII 2.7} 对 $K_\bullet$ 的类似结论. 设 $F$ 为凝聚 $\mathcal O_S$ 模, 其在 $K_\bullet(S)$ 中的类仍记作 $F$. 则 $g_*(F)$ 是复形 $Rg_*(F)=g_*(F)$ 在 $K_\bullet(X)$ 中的类, 亦即 $F$ 经 $g$ 看作 $\mathcal O_X$ 模所得的类. 为计算 $g^*(g_*(F))$, 须取
\[
 Lg^*(F)=F\otimes^L_{\mathcal O_X}\mathcal O_S.
\]
的类. 为此, 考虑 Koszul 复形
\[
\begin{aligned}
0&\longrightarrow \Lambda^r(E)\otimes_{\mathcal O_S}S(E)\longrightarrow
\Lambda^{r-1}(E)\otimes_{\mathcal O_S}S(E)\longrightarrow \cdots\\
&\longrightarrow E\otimes_{\mathcal O_S}S(E)\longrightarrow S(E)\longrightarrow 0,
\end{aligned}
\]
其中 $S(E)$ 为 $E$ 的对称代数. 这是 $\mathcal O_S$ 在 $S(E)$ 上的一个消解 (\hyperlink{sga6.e06.sec.1.10}{VI 1.11}). 相应的 $X$ 上的复形记作 $L_\bullet$, 其第 $k$ 项为 $\Lambda^k(f^*E)$, 因此是由有限型局部自由 $\mathcal O_X$ 模构成的 $\mathcal O_S$ 的消解. 于是
\[
 L_\bullet\otimes_{\mathcal O_X}F=\mathcal O_S\otimes^L_{\mathcal O_X}F,
\]
从而
\[
\begin{aligned}
 g^*(g_*(F))
 &=\cl(L_\bullet\otimes_{\mathcal O_X}F)
   =\sum_{i=0}^r(-1)^i\cl(L_i\otimes_{\mathcal O_X}F)\\
 &=\sum_{i=0}^r(-1)^i\cl(\Lambda^i(E)\otimes_{\mathcal O_S}F)
   =F\lambda_{-1}(E).
\end{aligned}
\]
由可加性得到命题.

\sgaStableTarget{sga6.e09.sec.2}\subsection*{2. 分裂环面下的主丛}

\sgaStableTarget{sga6.e09.para.2-1}\textbf{2.1.} 设 $S$ 为诺特概形, $T$ 为 $S$ 上的环面. 假定 $T$ 分裂; 因而有同构 $T\simeq D_S(M)$, 其中 $M$ 为秩 $r$ 的自由阿贝尔群, 或等价地 $T\simeq \mathbb G_{m,S}^r$. 设 $P$ 为 $T$ 下的挠子. 则 $P$ 在 $S$ 上仿射, 由一个 $M$ 型分次 $\mathcal O_S$ 代数 $A$ 定义. 若
\[
 A=\bigoplus_{m\in M}A_m,
\]
则对每个 $m$, $A_m$ 都是可逆 $\mathcal O_S$ 模, 而对 $m,m'\in M$, 由 $A$ 的乘法定义的同态
\[
 A_m\otimes_{\mathcal O_S}A_{m'}\longrightarrow A_{m+m'}
\]
为同构 (SGA 3 VIII 4.1). 因此定义乘法同态
\[
 \eta:M\longrightarrow K^\circ(S)
\]
将 $m\in M$ 对应于 $A_m$ 在 $K^\circ(S)$ 中的类.

\sgaStableTarget{sga6.e09.prop.2-2}\textbf{命题 2.2.} 沿用\hyperlink{sga6.e09.para.2-1}{2.1} 的记号与假设, 设 $f:P\to S$ 为 $P$ 的结构态射. 则同态
\[
 f^*:K_\bullet(S)\longrightarrow K_\bullet(P)
\]
为满射, 其核是 $K_\bullet(S)$ 中由各子群 $(\eta(m)-1)K_\bullet(S)$ 生成的子群, 其中 $m\in M$; 等价地, 只需令 $m$ 遍历 $M$ 的一组生成元.

先假定 $r=1$, 即 $T=\mathbb G_{m,S}$. 则 $M=\mathbb Z$, 且对每个 $m\in M$ 有
\[
 A_m\simeq A_1^{\otimes m}.
\]
若 $A_1=L$, 则 $A=\bigoplus_{n\in\mathbb Z}L^{\otimes n}$, 这说明 $P$ 可等同于丛 $X=\mathbb V(L)$ 中零截面的补集. 经零截面将 $S$ 看作 $X$ 的闭子概形, 由\hyperlink{sga6.e09.prop.1-1}{1.1}, 有正合列
\[
 K_\bullet(S)\xrightarrow{g_*}K_\bullet(X)\longrightarrow K_\bullet(P)\longrightarrow 0,
\]
其中 $g:S\to X$ 为 $X$ 的零截面. 此外, $f^*$ 分解为
\[
 K_\bullet(S)\xrightarrow{f'^*}K_\bullet(X)\longrightarrow K_\bullet(P),
\]
其中 $f'$ 为 $X$ 的结构态射. 由\hyperlink{sga6.e09.prop.1-6}{1.6}, $f'^*$ 是同构, 其逆为 $g^*$. 因而有正合列
\[
 K_\bullet(S)\xrightarrow{g^*g_*}K_\bullet(S)\xrightarrow{f^*}K_\bullet(P)\longrightarrow 0,
\]
由\hyperlink{sga6.e09.prop.1-8}{1.8}, 得到
\[
 K_\bullet(P)\simeq K_\bullet(S)/\lambda_{-1}(L)K_\bullet(S)=K_\bullet(S)/(1-L)K_\bullet(S),
\]
其中 $L$ 为 $L$ 在 $K^\circ(S)$ 中的类. 对每个 $n\in\mathbb Z$, $L^n-1$ 均为 $1-L$ 的倍数, 这就证明了 $P$ 的秩为一时的\hyperlink{sga6.e09.prop.2-2}{2.2}.

再对 $P$ 的秩 $r$ 作归纳. 假定结论对秩 $<r$ 的分裂环面下的挠子成立. 写成 $M=M_1\oplus M_2$, 其中 $M_1\simeq\mathbb Z^{r-1}$, $M_2\simeq\mathbb Z$. 置 $A'=\bigoplus_{m\in M_1}A_m$, $A''=\bigoplus_{m\in M_2}A_m$, 则 $A=A'\otimes_{\mathcal O_S}A''$.
设 $P_1=\Spec(A')$. 则 $P_1$ 为 $D_S(M_1)$ 下的挠子, 而 $P$ 在 $P_1$ 上是群 $\mathbb G_{m,P_1}$ 下的挠子, 其代数为 $A''$ 在 $P_1$ 上的逆像. 由上述秩一情形, $K_\bullet(P)$ 是 $K_\bullet(P_1)$ 模去由各 $(\eta(m_2)-1)K_\bullet(P_1)$ 生成的子群所得的商群, 其中 $m_2\in M_2$. 由归纳假设, $K_\bullet(P_1)$ 是 $K_\bullet(S)$ 模去由各 $(\eta(m_1)-1)K_\bullet(S)$ 生成的子群所得的商群, 其中 $m_1\in M_1$. 考虑关系
\[
 \eta(m_1+m_2)-1=\eta(m_1)\eta(m_2)-1
 =\eta(m_1)(\eta(m_2)-1)+\eta(m_1)-1,
\]
可知 $K_\bullet(P)$ 是 $K_\bullet(S)$ 模去由各 $(\eta(m)-1)K_\bullet(S)$ 生成的子群所得的商群, 其中 $m\in M$; 等价地, 只需令 $m$ 遍历 $M$ 的一组生成元.

\sgaStableTarget{sga6.e09.cor.2-3}\textbf{推论 2.3.} 沿用\hyperlink{sga6.e09.para.2-1}{2.1} 的记号与假设, 假定 $S$ 正则. 则有典范同构
\[
 K^\circ(P)\simeq K^\circ(S)/\bigl((\eta(m)-1)_{m\in M}\bigr),
\]
其中 $((\eta(m)-1)_{m\in M})$ 为 $K^\circ(S)$ 中由各 $\eta(m)-1$ 生成的理想, 这里 $m$ 遍历 $M$ 的一组生成元.

若 $S$ 正则, 则典范同态 $K^\circ(S)\to K_\bullet(S)$ 为同构 (\hyperlink{sga6.e04.para.2-5}{IV 2.5}). 同样, 此时 $P$ 也正则, 且 $K^\circ(P)\to K_\bullet(P)$ 为同构. 因此推论立即由 2.2 得到.

\sgaStableTarget{sga6.e09.sec.3}\subsection*{3. 射影丛与旗丛}

\sgaStableTarget{sga6.e09.prop.3-1}\textbf{命题 3.1.} 设 $S$ 为诺特概形, $E$ 为秩 $n$ 的局部自由 $\mathcal O_S$ 模, $\pi=(p_1,\ldots,p_k)$ 为满足 $\sum p_i=n$ 的整数序列, $X=D_\pi(E)$ 为 $E$ 的 $\pi$ 型旗概形. 则典范同态
\[
 K_\bullet(S)\otimes_{K^\circ(S)}K^\circ(X)\longrightarrow K_\bullet(X)
\]
为同构.

沿用\hyperlink{sga6.e06.sec.4.5}{VI 4.5} 的记号. 于是 $L_i$ 为 $E$ 在 $X$ 上的逆像的典范商, $F_j$ 为 $L_j\to L_{j-1}$ 的核, $\lambda_i^{(j)}$ 为 $\Lambda^i(F_j)$ 在 $K^\circ(X)$ 中的类. 则各 $\lambda_i^{(j)}$ 生成 $K^\circ(S)$ 代数 $K^\circ(X)$ (\hyperlink{sga6.e06.thm.4-6}{VI 4.6}). 同样, 对 $S$ 的每个点 $s$, 各 $\lambda_i^{(j)}$ 在 $K^\circ(X_s)$ 中的像生成 $K^\circ(k(s))=\mathbb Z$ 代数 $K^\circ(X_s)$. 由于 $X_s$ 是域上的旗概形, 它正则, 因而 $K^\circ(X_s)\to K_\bullet(X_s)$ 为同构. 从而 $K^\circ(X)\to K_\bullet(X_s)$ 为满射. 由\hyperlink{sga6.e09.cor.1-3}{1.3}, 这蕴含
\[
 K_\bullet(S)\otimes_{K^\circ(S)}K^\circ(X)\longrightarrow K_\bullet(X)
\]
为满射, 其中第二个因子为 $K^\circ(X)$.

为证单射性, 可先归结为完全旗丛 $D(E)$ 的情形. 事实上, 假定
\[
 K_\bullet(S)\otimes_{K^\circ(S)}K^\circ(D(E))\longrightarrow K_\bullet(D(E))
\]
为单射. 因为 $K_\bullet(S)\to K_\bullet(D(E))$ 经由 $K_\bullet(X)$ 分解, 对应的同态分解为
\[
\begin{aligned}
 &K_\bullet(S)\otimes_{K^\circ(S)}K^\circ(X)\otimes_{K^\circ(X)}K^\circ(D(E))\\
 &\qquad\longrightarrow K_\bullet(X)\otimes_{K^\circ(X)}K^\circ(D(E))\\
 &\qquad\longrightarrow K_\bullet(D(E)).
\end{aligned}
\]
第一个同态为单射. 由\hyperlink{sga6.e06.lem.4-3}{VI 4.3}、\hyperlink{sga6.e06.thm.4-6}{4.6} 及\hyperlink{sga6.e06.cor.4-7}{4.7}, $K^\circ(D(E))$ 在 $K^\circ(X)$ 上自由, 因此这蕴含
\[
 K_\bullet(S)\otimes_{K^\circ(S)}K^\circ(X)\to K_\bullet(X).
\]
为单射. 于是可假定 $X=D(E)$. 对 $E$ 的秩作归纳. $n=1$ 的情形平凡. 假定性质对 $n-1$ 成立. 引入 $P=\mathbb P(E)$, 设 $F$ 为典范同态 $p^*(E)\to\mathcal O_P(1)$ 的核, 其中 $p$ 为 $P$ 的结构态射. 则 $X$ 可等同于 $P$ 上 $F$ 的旗概形, 由归纳假设, 只需为 $P$ 证明\hyperlink{sga6.e09.prop.3-1}{3.1}. 考虑\hyperlink{sga6.e06.thm.1-1}{VI 1.1}, 所需证明的结论可表为以下推论.

\sgaStableTarget{sga6.e09.cor.3-2}\textbf{推论 3.2.} 设 $S$ 为诺特概形, $E$ 为秩 $n$ 的局部自由 $\mathcal O_S$ 模, $X=\mathbb P(E)$ 为相应的射影丛, $f$ 为 $X$ 的结构态射, $\xi$ 为层 $\mathcal O_X(1)$ 在 $K^\circ(X)$ 中的类. 则同态
\[
 \varphi:(K_\bullet(S))^n\longrightarrow K_\bullet(X)
\]
定义为
\[
 \varphi((x_i))=\sum_{i=0}^{n-1}\xi^{-i}f^*(x_i)
\]
是同构.

我们已知这个同态为满射, 因为各元素 $\xi^i$, $i=0,\ldots,-n+1$, 构成 $K^\circ(X)$ 在 $K^\circ(S)$ 上的一组基. 为证它是单射, 使用同态 $f_*$; 由于 $f$ 射影, 该同态有定义. 由\hyperlink{sga6.e04.eq.2-12-4}{IV 2.12.4},
\[
 f_*\left(\sum_{i=0}^{n-1}\xi^{-i}f^*(x_i)\right)
 =\sum_{i=0}^{n-1}x_i f_*(\xi^{-i})=x_0,
\]
最后的等式来自 $-n+1\leq i<0$ 时 $f_*(\xi^i)=0$ (\hyperlink{sga6.e06.sec.2.8}{VI 2.8}). 若 $\varphi((x_i))=0$, 则由上知 $x_0=0$. 对 $i=1,\ldots,n-1$ 依次将相同论证用于 $\xi^i\varphi((x_i))$, 归纳可得所有 $x_i$ 均为零. 因此 $\varphi$ 为单射.

\sgaStableTarget{sga6.e09.para.3-3}\textbf{3.3.} 仍设 $S$ 为诺特概形, $E$ 为秩 $n$ 的局部自由 $\mathcal O_S$ 模. 以 $V=\mathbb V(E)$ 表示与 $E$ 相伴且在 $E$ 上反变的向量丛 (SGA 3 I 4.6), 并视为 $S$ 上的群概形. 设 $P$ 为 $V$ 下的挠子; 我们要计算 $K_\bullet(P)$.

回顾如何将 $P$ 嵌入 $S$ 上的射影丛. $V$ 下的挠子对应于 $E$ 被 $\mathcal O_S$ 的扩张 $F$. 事实上, 若有扩张
\sgaStableTarget{sga6.e09.eq.3-3-1}\[
0\longrightarrow\mathcal O_S\xrightarrow{s}F\xrightarrow{u}E\longrightarrow 0,
\tag{3.3.1}
\]
则为得到 $V$ 下的挠子, 考虑由 $s$ 定义的同态 $\mathbb V(F)\to\mathbb V(\mathcal O_S)$ 以及由 $T=1$ 定义的 $\mathbb V(\mathcal O_S)=\Spec(\mathcal O_S[T])$ 的截面, 并取
\[
 P=\mathbb V(F)\times_{\mathbb V(\mathcal O_S)}S,
\]
便得到 $V$ 下的挠子, 其中 $V$ 在 $P$ 上的作用由它在 $\mathbb V(F)$ 上的作用诱导.

反之, 设 $P$ 为 $V$ 下的挠子. 用 $\underline P$ 表示从 $S$ 到 $P$ 的同态芽组成的层. 它是从 $S$ 到 $V$ 的同态芽层下的主齐性层, 后一个层可等同于 $E$ 的对偶 $\check E$. 按经典方法, 这给出 $\mathcal O_S$ 被 $\check E$ 的扩张, 再取对偶便得到\hyperlink{sga6.e09.eq.3-3-1}{(3.3.1)}.

因此, 设 $P$ 为 $V$ 下的挠子, 对应于 $E$ 被 $\mathcal O_S$ 的扩张 $F$. $P$ 在 $\mathcal O_S$ 上的代数 $A(P)$ 为
\[
 S(F)\otimes_{S(\mathcal O_S)}\mathcal O_S,
\]
即
\[
 S(F)/(s-1)S(F),
\]
它等同于次数 $0$ 的齐次分式环 $S(F)_{(s)}$. 因此 $P$ 等同于 $\mathbb P(F)$ 的开集 $D_+(s)$. $P$ 在 $\mathbb P(F)$ 中的补集为 $\operatorname{Proj}(S(F)/sS(F))$, 由\hyperlink{sga6.e09.eq.3-3-1}{(3.3.1)} 可将其等同于 $\mathbb P(E)$. 因而有浸入
\[
\begin{tikzcd}[column sep=large,row sep=small]
\mathbb P(E) \arrow[r,hook] \arrow[dr] & \mathbb P(F) & P \arrow[l,hook'] \arrow[dl]\\
& S &
\end{tikzcd}
\]
其中 $\mathbb P(E)$ 在 $\mathbb P(F)$ 中闭, 而 $P$ 为 $\mathbb P(E)$ 的开补集.

\sgaStableTarget{sga6.e09.prop.3-4}\textbf{命题 3.4.} 设 $S$ 为诺特概形, $E$ 为秩 $n$ 的局部自由 $\mathcal O_S$ 模, $V=\mathbb V(E)$ 为相应的向量丛, $P$ 为 $V$ 下的挠子, $f$ 为 $P$ 的结构态射. 则典范同态
\[
 f^*:K_\bullet(S)\longrightarrow K_\bullet(P)
\]
为同构.

通过\hyperlink{sga6.e09.para.3-3}{3.3} 中的等同, 有正合列
\sgaStableTarget{sga6.e09.eq.3-4-1}\[
 K_\bullet(\mathbb P(E))\xrightarrow{i_*}K_\bullet(\mathbb P(F))\longrightarrow K_\bullet(P)\longrightarrow 0,
\tag{3.4.1}
\]
其中 $i$ 表示由 $u$ 定义的闭浸入 $\mathbb P(E)\to\mathbb P(F)$. 由\hyperlink{sga6.e09.cor.3-2}{3.2},
\[
 K_\bullet(\mathbb P(E))=K_\bullet(S)\otimes_{K^\circ(S)}K^\circ(\mathbb P(E)),
\qquad
 K_\bullet(\mathbb P(F))=K_\bullet(S)\otimes_{K^\circ(S)}K^\circ(\mathbb P(F)).
\]
由于 $\mathbb P(E)$ 和 $\mathbb P(F)$ 光滑, $i$ 为正则浸入, 因而有同态
\[
 i_*:K^\circ(\mathbb P(E))\longrightarrow K^\circ(\mathbb P(F)).
\]
$K_\bullet$ 群间的同态可由 $K^\circ$ 群间的同态与 $K_\bullet(S)$ 作张量积得到. 为此, 只须证明对 $a\in K_\bullet(S)$ 和 $b\in K^\circ(\mathbb P(E))$ 有
\[
 i_*(b\,h^*(a))=i_*(b)g^*(a),
\]
其中 $g$ 和 $h$ 分别为 $\mathbb P(F)$ 和 $\mathbb P(E)$ 的结构态射. 因为 $h=gi$, 由\hyperlink{sga6.e04.eq.2-12-4}{IV 2.12.4} 有
\[
 i_*(b\,h^*(a))=i_*(b\,i^*g^*(a))=g^*(a)i_*(b)
\]
%

故由\hyperlink{sga6.e09.eq.3-4-1}{(3.4.1)} 得到
\[
 K_\bullet(P)\simeq K_\bullet(S)\otimes_{K^\circ(S)}
 \bigl(K^\circ(\mathbb P(F))/\operatorname{Im}K^\circ(\mathbb P(E))\bigr).
\]
设 $\xi_F$ 和 $\xi_E$ 分别为 $\mathbb P(F)$ 和 $\mathbb P(E)$ 上典范可逆层的类. 由\hyperlink{sga6.e06.thm.1-1}{VI 1.1}, $K^\circ(\mathbb P(F))$ (分别为 $K^\circ(\mathbb P(E))$) 在 $K^\circ(S)$ 上以 $\xi_F^k$, $k=0,\ldots,-n$, (分别为 $\xi_E^k$, $k=0,\ldots,-n+1$) 为基. 投影公式给出
\[
 i_*(\xi_E^k)=i_*(i^*(\xi_F^k))=\xi_F^k i_*(1).
\]
若 $I$ 为 $\mathbb P(E)$ 在 $\mathbb P(F)$ 中的理想, 则 $I$ 可逆且 $i_*(1)=1-I$. 此外, $I$ 由 $S(F)$ 的理想 $J$ 定义, 它是 $S(F)\to S(E)$ 的核, 由 $F$ 中同构于 $\mathcal O_S$ 的子模生成. 因而
\[
 I\simeq \mathcal O_{\mathbb P(F)}(-1).
\]
故 $K^\circ(\mathbb P(E))$ 的像为各 $\xi_F^k-\xi_F^{k-1}$, $k=0,\ldots,-n+1$, 所生成的 $K^\circ(S)$ 子模. 因此
\[
 K^\circ(S)\longrightarrow K^\circ(\mathbb P(F))/\operatorname{Im}K^\circ(\mathbb P(E))
\]
为同构, 从而 $K_\bullet(S)\to K_\bullet(P)$ 也为同构.

\sgaStableTarget{sga6.e09.sec.4}\subsection*{4. 群 $\mathrm{GL}(n)_S$ 下的主丛}

\sgaStableTarget{sga6.e09.prop.4-1}\textbf{命题 4.1.} 设 $S$ 为诺特概形, $G=\mathrm{GL}(n)_S$, $P$ 为 $G$ 下的挠子, $E$ 为与 $P$ 对应的秩 $n$ 局部自由 $\mathcal O_S$ 模. 用 $E$ 也表示 $E$ 在 $K^\circ(S)$ 中的类, 设 $J$ 为 $E-n$ 生成的 $\lambda$ 理想. 若 $f$ 为 $P$ 的结构态射, 则 $f^*$ 因 $f$ 平坦而有定义, 是满射且定义出同构
\[
 K_\bullet(S)/J K_\bullet(S)\xrightarrow{\sim}K_\bullet(P).
\]

设 $E_0=(\mathcal O_S)^n$, 带有显然的基 $e_1,\ldots,e_n$, 并设 $E_0^{(i)}$ 为 $E_0$ 中由 $e_1,\ldots,e_i$ 生成的子模. 将每个 $S$ 概形 $S'$ 对应于 $E_{0,S'}$ 的对所有 $i$ 保持子模 $E_{0,S'}^{(i)}$ 不变的自同构所成的子群, 这个 $G$ 的子函子由 $G$ 的群子概形 $B$ 表示, 它是 $G$ 的 Borel 子群. fpqc 拓扑下的商层 $G/B$ 由 $E_0$ 的旗概形表示, 记作 $X_0$. 设 $L_0^{(i)}$ 为 $E_0$ 在 $X_0$ 上的逆像的典范商, $G_0^{(i)}$ 为 $E_{0,X_0}\to L_0^{(i)}$ 的核, 并置
\[
 F_0^{(i)}=G_0^{(i-1)}/G_0^{(i)}\simeq \ker(L_0^{(i)}\to L_0^{(i-1)}).
\]

% SGA-SRC-001466
此外, $B$ 为其幺幂根基 $U$ 与保持各 $e_i$ 不变的极大子环面 $T$ 的半直积. 设 $Y_0=G/U$. 则 $Y_0$ 是 $X_0$ 上群 $T_{X_0}$ 下的挠子. 若 $S'$ 为 $S$ 概形, 给定 $Y_0$ 的一个 $S'$ 点等价于在某个忠实平坦且拟紧的 $S'$ 概形 $S''$ 上给定一个 $E_{0,S''}$ 的自同构模去 $U_{S''}$ 的幺幂元素; 亦即给定由自由子模构成的 $E_{0,S''}$ 的旗, 并对每个 $i$ 给定一类在 $\mathcal O_{S''}$ 上自由的元素 $f_i^\alpha$, 使得 $f_i^\alpha-f_i^\beta\in G_{0,S''}^{(i)}$. 等价地, 给定 $E_{0,S''}$ 的旗以及对每个 $i$ 给定 $F_{0,S''}^{(i)}$ 的一组基, 这些数据须满足 $S''\times_{S'}S''$ 上的黏合条件. 由 fpqc 下降, 这等价于在 $S'$ 上给定同样的数据. 因此 $Y_0$ 为代数
\[
 A_0=\bigoplus_{k_i\in\mathbb Z}(F_0^{(1)})^{\otimes k_1}\otimes\cdots\otimes(F_0^{(n)})^{\otimes k_n}.
\]

定义的仿射 $X_0$ 概形. 群 $G$ 作用在 $E_0$、带有各层 $L_0^{(i)}$ 的 $X_0$ 以及 $Y_0$ 上. 将这些对象与 $P$ 作乘积后模去 $G$, 得到 $P$ 所伴的层 $E$、$E$ 的旗概形 $X=D(E)=P/B$ 以及概形 $Y=P/U$. $X$ 上的典范层仍用 $X_0$ 上相应层的字母表示. 于是有图表
\[
\begin{array}{ccccc}
&& Y=P/U &&\\[-0.2em]
& \nearrow && \searrow &\\[-0.2em]
P &&&& X=P/B=D(E)\\[-0.2em]
& \searrow && \swarrow &\\[-0.2em]
&& S &&
\end{array}
\]
注意 $P$ 是 $Y$ 上结构群为 $U_Y$ 的挠子, 而 $Y$ 是 $X$ 上结构群为 $T_X$ 的挠子, 由 $\mathcal O_X$ 代数
\[
 A=\bigoplus_{k_i\in\mathbb Z}(F^{(1)})^{\otimes k_1}\otimes\cdots\otimes(F^{(n)})^{\otimes k_n}.
\]

定义. 因此同态 $K_\bullet(S)\to K_\bullet(P)$ 分解为
\[
 K_\bullet(S)\longrightarrow K_\bullet(X)\longrightarrow K_\bullet(Y)\longrightarrow K_\bullet(P).
\]
由\hyperlink{sga6.e09.prop.3-1}{3.1}, $K_\bullet(X)=K_\bullet(S)\otimes_{K^\circ(S)}K^\circ(X)$. 设 $\xi_i$ 为 $F^{(i)}$ 在 $K^\circ(X)$ 中的类. 由\hyperlink{sga6.e09.prop.2-2}{2.2}, $K_\bullet(Y)$ 是 $K_\bullet(X)$ 模去由各 $(\xi_i-1)K_\bullet(X)$ 生成的子模所得的商. 若 $I$ 为 $K^\circ(X)$ 中由各 $\xi_i-1$ 生成的理想, 则
\[
 K_\bullet(Y)=K_\bullet(S)\otimes_{K^\circ(S)}(K^\circ(X)/I).
\]
由\hyperlink{sga6.e06.cor.4-8}{VI 4.8}, $K^\circ(X)$ 为各 $\xi_i-1$ 生成的 $K^\circ(S)$ 代数, 其关系由
\[
 \gamma_t(E-n)=\prod_{i=1}^n\gamma_t(\xi_i-1).
\]
导出. 因而 $K^\circ(X)/I=K^\circ(S)/J$, 其中 $J$ 为各 $\gamma^j(E-n)$, $j=1,\ldots,n$, 所生成的理想, 因此是 $E-n$ 生成的 $\lambda$ 理想. 故
\[
 K_\bullet(Y)\simeq K_\bullet(S)/J K_\bullet(S).
\]
下述引理完成证明.

\sgaStableTarget{sga6.e09.lem.4-2}\textbf{引理 4.2.} 设 $S$ 为诺特概形, $B$ 为约化 $S$ 群的抛物子群, $P$ 为 $B$ 的幺幂根基 $U$ 下的挠子. 则有同构
\[
 K_\bullet(S)\xrightarrow{\sim}K_\bullet(P).
\]

事实上, 存在 $U$ 的一列群子概形
\[
 U_0=U\supset U_1\supset U_2\supset\cdots\supset U_p=e,
\]
它们在 $B$ 中闭, 因 $B$ 诺特而为有限列, 且对每个 $i$, $U_i/U_{i+1}$ 都同构于一个 $S$ 向量群 (SGA 3 XXVI 2.1). 因而 $P\to S$ 分解为一串向量群下的挠子, 结论由 3.4 得到.

\sgaStableTarget{sga6.e09.cor.4-3}\textbf{推论 4.3.} 设 $S$ 为诺特概形, $G=\mathrm{GL}(n)_S$. 则典范同态
\[
 K_\bullet(S)\longrightarrow K_\bullet(G)
\]
为同构.

这立即由\hyperlink{sga6.e09.prop.4-1}{4.1} 应用于 $P=G$ 得到, 因为此时 $E=\mathcal O_S^n$.

\sgaStableTarget{sga6.e09.cor.4-4}\textbf{推论 4.4.} 在\hyperlink{sga6.e09.prop.4-1}{4.1} 的假设下, 若 $S$ 正则, 则有同构
\[
 K^\circ(S)/J\xrightarrow{\sim}K^\circ(P).
\]

若 $S$ 正则, 则 $G$ 以及 $P$ 也正则. 因此 $K^\circ(S)\simeq K_\bullet(S)$, $K^\circ(P)\simeq K_\bullet(P)$, 结论由 4.1 得到.

